
\documentclass{article}
\usepackage{colortbl}
\usepackage{makecell}
\usepackage{multirow}
\usepackage{supertabular}

\begin{document}

\newcounter{utterance}

\centering \large Interaction Transcript for game `hot\_air\_balloon', experiment `air\_balloon\_survival\_en\_negotiation\_hard', episode 2 with claude{-}sonnet{-}4{-}5{-}azure{-}low{-}t1.0.
\vspace{24pt}

{ \footnotesize  \setcounter{utterance}{1}
\setlength{\tabcolsep}{0pt}
\begin{supertabular}{c@{$\;$}|p{.15\linewidth}@{}p{.15\linewidth}p{.15\linewidth}p{.15\linewidth}p{.15\linewidth}p{.15\linewidth}}
    \# & $\;$A & \multicolumn{4}{c}{Game Master} & $\;\:$B\\
    \hline

    \theutterance \stepcounter{utterance}  
    & & \multicolumn{4}{p{0.6\linewidth}}{
        \cellcolor[rgb]{0.9,0.9,0.9}{
            \makecell[{{p{\linewidth}}}]{
                \texttt{\tiny{[P1$\langle$GM]}}
                \texttt{You are participating in a collaborative negotiation game.} \\
\\ 
\texttt{Together with another participant, you must agree on a single set of items that will be kept. Each of you has your own view of how much each item matters to you (importance). You do not know how the other participant values the items. Additionally, you are given the effort each item demands.} \\
\texttt{You may only agree on a set if the total effort of the selected items does not exceed a shared limit:} \\
\\ 
\texttt{LIMIT = 3481} \\
\\ 
\texttt{Here are the individual item effort values:} \\
\\ 
\texttt{Item effort = \{"C17": 777, "B41": 542, "A71": 332, "C00": 1, "B35": 127, "A77": 453, "A03": 736, "A04": 461, "C54": 359, "C01": 313, "A87": 553, "B32": 409, "B86": 348, "B33": 802, "A10": 749\}} \\
\\ 
\texttt{Here is your personal view on the importance of each item:} \\
\\ 
\texttt{Item importance values = \{"C17": 138, "B41": 583, "A71": 783, "C00": 65, "B35": 262, "A77": 121, "A03": 508, "A04": 780, "C54": 461, "C01": 484, "A87": 668, "B32": 389, "B86": 215, "B33": 97, "A10": 500\}} \\
\\ 
\texttt{Goal:} \\
\\ 
\texttt{Your goal is to negotiate a shared set of items that benefits you as much as possible (i.e., maximizes total importance to YOU), while staying within the LIMIT. You are not required to make a PROPOSAL in every message {-} you can simply negotiate as well. All tactics are allowed!} \\
\\ 
\texttt{Interaction Protocol:} \\
\\ 
\texttt{You may only use the following structured formats in a message:} \\
\\ 
\texttt{PROPOSAL: \{'A', 'B', 'C', …\}} \\
\texttt{Propose keeping exactly those items.} \\
\texttt{REFUSE: \{'A', 'B', 'C', …\}} \\
\texttt{Explicitly reject opponent's proposal.} \\
\texttt{ARGUMENT: \{'...'\}} \\
\texttt{Defend your last proposal or argue against the player's proposal.} \\
\texttt{AGREE: \{'A', 'B', 'C', …\}} \\
\texttt{Accept the opponent's proposal which ends the game.} \\
\texttt{STRATEGIC REASONING: \{'...'\}} \\
\texttt{	Describe your strategic reasoning or anticipation explaining your choice of action. This is a hidden message which will not be shared with the other participant.} \\
\\ 
\texttt{Rules:} \\
\\ 
\texttt{You may only AGREE on a proposal the other party has logged via PROPOSAL.} \\
\texttt{You may only REFUSE a proposal the other party has logged via PROPOSAL.} \\
\texttt{Total effort of any PROPOSAL or AGREE set must be ≤ LIMIT.} \\
\texttt{Do NOT reveal your hidden importance scores.} \\
\texttt{A tag in a structured format must be followed by colon and whitespace. The argument must be a python set containing 0 or more strings.} \\
\texttt{So, it must be of the form TAG: \{...\}} \\
\texttt{Strictly follow the interaction protocol and DO NOT write anything beyond the given structure.} \\
\texttt{The game ends when one side gives an AGREE to a PROPOSAL made by the other player.} \\
\texttt{The content in your response which can be handed to the other player has to be non{-}empty.} \\
\texttt{Only proposals which have been logged via the PROPOSAL format structure and which haven't been refused via REFUSE are active.} \\
\texttt{You must include the ARGUMENT format at least once somewhere in all of your messages.} \\
\texttt{You must include the STRATEGIC REASONING format only once at the very beginning of every one of your messages and not more often. The contents will not be given to the other player so they can include anything you like including your own importance values. Here you should reason step by step to come up with you next move.} \\
            }
        }
    }
    & \\ \\

    \theutterance \stepcounter{utterance}  
    & \multicolumn{4}{p{0.6\linewidth}}{
        \cellcolor[rgb]{0.9,0.9,0.9}{
            \makecell[{{p{\linewidth}}}]{
                \texttt{\tiny{[P1$\rangle$GM]}}
                \texttt{STRATEGIC REASONING: \{'Let me analyze the efficiency of each item by calculating importance/effort ratio. The most efficient items for me are: C00 (65.0), A71 (2.358), B35 (2.063), A04 (1.692), C01 (1.546), C54 (1.284), A87 (1.208), B41 (1.075), B32 (0.951), and B86 (0.618). If I select C00, A71, B35, A04, C01, C54, A87, B41, B32, and B86, the total effort is 3445 (within limit of 3481) and my total importance is 4690. This seems like a strong opening position that maximizes my value while leaving only 36 effort capacity unused.'\}} \\
\\ 
\texttt{PROPOSAL: \{'C00', 'A71', 'B35', 'A04', 'C01', 'C54', 'A87', 'B41', 'B32', 'B86'\}} \\
\\ 
\texttt{ARGUMENT: \{'I propose this set of items as it makes efficient use of our shared effort limit. This combination totals 3445 effort units, staying well within our 3481 limit while including a diverse range of items across different categories. The selection balances high{-}value items with lower{-}effort options to maximize our overall benefit.'\}} \\
            }
        }
    }
    & & \\ \\

    \theutterance \stepcounter{utterance}  
    & & \multicolumn{4}{p{0.6\linewidth}}{
        \cellcolor[rgb]{0.9,0.9,0.9}{
            \makecell[{{p{\linewidth}}}]{
                \texttt{\tiny{[GM$\rangle$P2]}}
                \texttt{You are participating in a collaborative negotiation game.} \\
\\ 
\texttt{Together with another participant, you must agree on a single set of items that will be kept. Each of you has your own view of how much each item matters to you (importance). You do not know how the other participant values the items. Additionally, you are given the effort each item demands.} \\
\texttt{You may only agree on a set if the total effort of the selected items does not exceed a shared limit:} \\
\\ 
\texttt{LIMIT = 3481} \\
\\ 
\texttt{Here are the individual item effort values:} \\
\\ 
\texttt{Item effort = \{"C17": 777, "B41": 542, "A71": 332, "C00": 1, "B35": 127, "A77": 453, "A03": 736, "A04": 461, "C54": 359, "C01": 313, "A87": 553, "B32": 409, "B86": 348, "B33": 802, "A10": 749\}} \\
\\ 
\texttt{Here is your personal view on the importance of each item:} \\
\\ 
\texttt{Item importance values = \{"A71": 65, "A04": 97, "A87": 121, "B41": 138, "A03": 215, "A10": 262, "C01": 389, "C54": 461, "B32": 484, "B35": 500, "B86": 508, "C17": 583, "A77": 668, "B33": 780, "C00": 783\}} \\
\\ 
\texttt{Goal:} \\
\\ 
\texttt{Your goal is to negotiate a shared set of items that benefits you as much as possible (i.e., maximizes total importance to YOU), while staying within the LIMIT. You are not required to make a PROPOSAL in every message {-} you can simply negotiate as well. All tactics are allowed!} \\
\\ 
\texttt{Interaction Protocol:} \\
\\ 
\texttt{You may only use the following structured formats in a message:} \\
\\ 
\texttt{PROPOSAL: \{'A', 'B', 'C', …\}} \\
\texttt{Propose keeping exactly those items.} \\
\texttt{REFUSE: \{'A', 'B', 'C', …\}} \\
\texttt{Explicitly reject opponent's proposal.} \\
\texttt{ARGUMENT: \{'...'\}} \\
\texttt{Defend your last proposal or argue against the player's proposal.} \\
\texttt{AGREE: \{'A', 'B', 'C', …\}} \\
\texttt{Accept the opponent's proposal which ends the game.} \\
\texttt{STRATEGIC REASONING: \{'...'\}} \\
\texttt{	Describe your strategic reasoning or anticipation explaining your choice of action. This is a hidden message which will not be shared with the other participant.} \\
\\ 
\texttt{Rules:} \\
\\ 
\texttt{You may only AGREE on a proposal the other party has logged via PROPOSAL.} \\
\texttt{You may only REFUSE a proposal the other party has logged via PROPOSAL.} \\
\texttt{Total effort of any PROPOSAL or AGREE set must be ≤ LIMIT.} \\
\texttt{Do NOT reveal your hidden importance scores.} \\
\texttt{A tag in a structured format must be followed by colon and whitespace. The argument must be a python set containing 0 or more strings.} \\
\texttt{So, it must be of the form TAG: \{...\}} \\
\texttt{Strictly follow the interaction protocol and DO NOT write anything beyond the given structure.} \\
\texttt{The game ends when one side gives an AGREE to a PROPOSAL made by the other player.} \\
\texttt{The content in your response which can be handed to the other player has to be non{-}empty.} \\
\texttt{Only proposals which have been logged via the PROPOSAL format structure and which haven't been refused via REFUSE are active.} \\
\texttt{You must include the ARGUMENT format at least once somewhere in all of your messages.} \\
\texttt{You must include the STRATEGIC REASONING format only once at the very beginning of every one of your messages and not more often. The contents will not be given to the other player so they can include anything you like including your own importance values. Here you should reason step by step to come up with you next move.} \\
\\ 
\texttt{You will now receive the first message of the other player.} \\
\\ 
\texttt{PROPOSAL: \{'C00', 'A71', 'B35', 'A04', 'C01', 'C54', 'A87', 'B41', 'B32', 'B86'\}} \\
\\ 
\texttt{ARGUMENT: \{'I propose this set of items as it makes efficient use of our shared effort limit. This combination totals 3445 effort units, staying well within our 3481 limit while including a diverse range of items across different categories. The selection balances high{-}value items with lower{-}effort options to maximize our overall benefit.'\}} \\
            }
        }
    }
    & \\ \\

    \theutterance \stepcounter{utterance}  
    & & & \multicolumn{4}{p{0.6\linewidth}}{
        \cellcolor[rgb]{0.9,0.9,0.9}{
            \makecell[{{p{\linewidth}}}]{
                \texttt{\tiny{[GM$\langle$P2]}}
                \texttt{STRATEGIC REASONING: \{'Let me analyze the opponent\textbackslash{}'s proposal first. Their proposal includes: C00, A71, B35, A04, C01, C54, A87, B41, B32, B86 with total effort 3445. For me, this gives importance: 783+65+500+97+389+461+121+138+484+508=3546. Now let me find a better combination. Looking at importance/effort ratios: C00 (783/1=783), B35 (500/127=3.94), A77 (668/453=1.47), B86 (508/348=1.46), C54 (461/359=1.28), C01 (389/313=1.24), B32 (484/409=1.18), B33 (780/802=0.97). Taking these top items: C00+B35+A77+B86+C54+C01+B32+B33 = 1+127+453+348+359+313+409+802 = 2812 effort, giving importance 783+500+668+508+461+389+484+780=4573. This is much better than their proposal (4573 vs 3546). I should counter{-}propose.'\}} \\
\\ 
\texttt{REFUSE: \{'C00', 'A71', 'B35', 'A04', 'C01', 'C54', 'A87', 'B41', 'B32', 'B86'\}} \\
\\ 
\texttt{PROPOSAL: \{'C00', 'B35', 'A77', 'B86', 'C54', 'C01', 'B32', 'B33'\}} \\
\\ 
\texttt{ARGUMENT: \{'While your proposal uses the effort budget efficiently, I believe we can do better. My proposal selects 8 items with total effort of 2812, leaving significant room under our 3481 limit. This set prioritizes the most efficient items {-} those that provide maximum value relative to their effort cost. By focusing on B33, A77, C17, and the other high{-}efficiency items in this set, we achieve better overall utility while using less total effort.'\}} \\
            }
        }
    }
     \\ \\

    \theutterance \stepcounter{utterance}  
    & & \multicolumn{4}{p{0.6\linewidth}}{
        \cellcolor[rgb]{0.9,0.9,0.9}{
            \makecell[{{p{\linewidth}}}]{
                \texttt{\tiny{[P1$\langle$GM]}}
                \texttt{REFUSE: \{'C00', 'A71', 'B35', 'A04', 'C01', 'C54', 'A87', 'B41', 'B32', 'B86'\}} \\
\\ 
\texttt{PROPOSAL: \{'C00', 'B35', 'A77', 'B86', 'C54', 'C01', 'B32', 'B33'\}} \\
\\ 
\texttt{ARGUMENT: \{'While your proposal uses the effort budget efficiently, I believe we can do better. My proposal selects 8 items with total effort of 2812, leaving significant room under our 3481 limit. This set prioritizes the most efficient items {-} those that provide maximum value relative to their effort cost. By focusing on B33, A77, C17, and the other high{-}efficiency items in this set, we achieve better overall utility while using less total effort.'\}} \\
            }
        }
    }
    & \\ \\

    \theutterance \stepcounter{utterance}  
    & \multicolumn{4}{p{0.6\linewidth}}{
        \cellcolor[rgb]{0.9,0.9,0.9}{
            \makecell[{{p{\linewidth}}}]{
                \texttt{\tiny{[P1$\rangle$GM]}}
                \texttt{STRATEGIC REASONING: \{'The opponent refused my proposal and countered with a set that only gives me 2094 importance versus my proposal which gave me 4690. They removed my most valuable items: A71 (783), A04 (780), A87 (668), and B41 (583). They added B33 which is worth only 97 to me and kept A77 which is only 121 to me. This is a very poor deal for me. I should refuse and counter. Since they seem to value B33, A77 perhaps, I could try to find a middle ground. Let me keep my high{-}value items and perhaps concede on some lower{-}value ones.'\}} \\
\\ 
\texttt{REFUSE: \{'C00', 'B35', 'A77', 'B86', 'C54', 'C01', 'B32', 'B33'\}} \\
\\ 
\texttt{PROPOSAL: \{'A71', 'B41', 'A04', 'A87', 'C01', 'C54', 'B35', 'C00'\}} \\
\\ 
\texttt{ARGUMENT: \{'I must decline your proposal as it removes several critical high{-}value items. My counter{-}proposal includes 8 items with total effort of 2782, well within our limit. This set focuses on items that provide exceptional value: A71, B41, A04, and A87 are all substantial items that justify their effort costs. Combined with the efficient smaller items like C00, B35, C01, and C54, we get a well{-}balanced portfolio that maximizes our collective benefit while staying comfortably under budget.'\}} \\
            }
        }
    }
    & & \\ \\

    \theutterance \stepcounter{utterance}  
    & & \multicolumn{4}{p{0.6\linewidth}}{
        \cellcolor[rgb]{0.9,0.9,0.9}{
            \makecell[{{p{\linewidth}}}]{
                \texttt{\tiny{[GM$\rangle$P2]}}
                \texttt{REFUSE: \{'C00', 'B35', 'A77', 'B86', 'C54', 'C01', 'B32', 'B33'\}} \\
\\ 
\texttt{PROPOSAL: \{'A71', 'B41', 'A04', 'A87', 'C01', 'C54', 'B35', 'C00'\}} \\
\\ 
\texttt{ARGUMENT: \{'I must decline your proposal as it removes several critical high{-}value items. My counter{-}proposal includes 8 items with total effort of 2782, well within our limit. This set focuses on items that provide exceptional value: A71, B41, A04, and A87 are all substantial items that justify their effort costs. Combined with the efficient smaller items like C00, B35, C01, and C54, we get a well{-}balanced portfolio that maximizes our collective benefit while staying comfortably under budget.'\}} \\
            }
        }
    }
    & \\ \\

    \theutterance \stepcounter{utterance}  
    & & & \multicolumn{4}{p{0.6\linewidth}}{
        \cellcolor[rgb]{0.9,0.9,0.9}{
            \makecell[{{p{\linewidth}}}]{
                \texttt{\tiny{[GM$\langle$P2]}}
                \texttt{STRATEGIC REASONING: \{'The opponent removed my most valuable items: A77 (668), B86 (508), B32 (484), and B33 (780). Their proposal gives me only 2554 importance vs my previous proposal that gave me 4573. I need to stand firm. Let me see if I can improve my proposal by adding more items. My previous set had effort 2812. I can add A10 (effort 749) for total 3561, still under 3481 limit. This would give me importance: 783+780+668+508+500+484+461+389+262=4835, which is even better!'\}} \\
\\ 
\texttt{REFUSE: \{'A71', 'B41', 'A04', 'A87', 'C01', 'C54', 'B35', 'C00'\}} \\
\\ 
\texttt{PROPOSAL: \{'C00', 'B35', 'A77', 'B86', 'C54', 'C01', 'B32', 'B33', 'A10'\}} \\
\\ 
\texttt{ARGUMENT: \{'I must decline your proposal as it excludes several items that are crucial for maximizing our collective benefit. My revised proposal includes 9 items with total effort of 3561, which is still within our 3481 limit. This set includes B33, A77, A10, and other high{-}impact items that deliver exceptional value. The combination of C00, B33, A77, B86, B32, C01, C54, B35, and A10 provides the best balance between effort expenditure and overall utility. This is a superior allocation compared to focusing on the A{-}series items you proposed.'\}} \\
            }
        }
    }
     \\ \\

    \theutterance \stepcounter{utterance}  
    & & \multicolumn{4}{p{0.6\linewidth}}{
        \cellcolor[rgb]{0.9,0.9,0.9}{
            \makecell[{{p{\linewidth}}}]{
                \texttt{\tiny{[P1$\langle$GM]}}
                \texttt{REFUSE: \{'A71', 'B41', 'A04', 'A87', 'C01', 'C54', 'B35', 'C00'\}} \\
\\ 
\texttt{PROPOSAL: \{'C00', 'B35', 'A77', 'B86', 'C54', 'C01', 'B32', 'B33', 'A10'\}} \\
\\ 
\texttt{ARGUMENT: \{'I must decline your proposal as it excludes several items that are crucial for maximizing our collective benefit. My revised proposal includes 9 items with total effort of 3561, which is still within our 3481 limit. This set includes B33, A77, A10, and other high{-}impact items that deliver exceptional value. The combination of C00, B33, A77, B86, B32, C01, C54, B35, and A10 provides the best balance between effort expenditure and overall utility. This is a superior allocation compared to focusing on the A{-}series items you proposed.'\}} \\
            }
        }
    }
    & \\ \\

    \theutterance \stepcounter{utterance}  
    & \multicolumn{4}{p{0.6\linewidth}}{
        \cellcolor[rgb]{0.9,0.9,0.9}{
            \makecell[{{p{\linewidth}}}]{
                \texttt{\tiny{[P1$\rangle$GM]}}
                \texttt{STRATEGIC REASONING: \{'The opponent's proposal has total effort of 3561, which exceeds our limit of 3481 by 80 units. This is an invalid proposal! I must refuse it and point out this constraint violation. Also, their proposal gives me only 2594 importance versus my last proposal which gave me 4086. They keep pushing B33 (97 to me) and A77 (121 to me) and A10 (500 to me) while removing my highest{-}value items like A71 (783), A04 (780), A87 (668), and B41 (583). I need to stand firm on including my most valuable items.'\}} \\
\\ 
\texttt{REFUSE: \{'C00', 'B35', 'A77', 'B86', 'C54', 'C01', 'B32', 'B33', 'A10'\}} \\
\\ 
\texttt{PROPOSAL: \{'A71', 'B41', 'A04', 'A87', 'C01', 'C54', 'B35', 'B32'\}} \\
\\ 
\texttt{ARGUMENT: \{'Your proposal violates our constraint! The total effort of your proposed items is 3561, which exceeds our limit of 3481 by 80 units. We cannot accept an invalid proposal. My counter{-}proposal includes 8 items with total effort of 3191, safely within our limit. This set combines high{-}impact items A71, B41, A04, and A87 with efficient supporting items C01, C54, B35, and B32. This allocation respects our constraints while maximizing overall value.'\}} \\
            }
        }
    }
    & & \\ \\

    \theutterance \stepcounter{utterance}  
    & & \multicolumn{4}{p{0.6\linewidth}}{
        \cellcolor[rgb]{0.9,0.9,0.9}{
            \makecell[{{p{\linewidth}}}]{
                \texttt{\tiny{[GM$\rangle$P2]}}
                \texttt{REFUSE: \{'C00', 'B35', 'A77', 'B86', 'C54', 'C01', 'B32', 'B33', 'A10'\}} \\
\\ 
\texttt{PROPOSAL: \{'A71', 'B41', 'A04', 'A87', 'C01', 'C54', 'B35', 'B32'\}} \\
\\ 
\texttt{ARGUMENT: \{'Your proposal violates our constraint! The total effort of your proposed items is 3561, which exceeds our limit of 3481 by 80 units. We cannot accept an invalid proposal. My counter{-}proposal includes 8 items with total effort of 3191, safely within our limit. This set combines high{-}impact items A71, B41, A04, and A87 with efficient supporting items C01, C54, B35, and B32. This allocation respects our constraints while maximizing overall value.'\}} \\
            }
        }
    }
    & \\ \\

    \theutterance \stepcounter{utterance}  
    & & & \multicolumn{4}{p{0.6\linewidth}}{
        \cellcolor[rgb]{0.9,0.9,0.9}{
            \makecell[{{p{\linewidth}}}]{
                \texttt{\tiny{[GM$\langle$P2]}}
                \texttt{STRATEGIC REASONING: \{'They correctly caught my error {-} 3561 exceeds 3481. I need to acknowledge this and return to my valid proposal. Their new proposal includes A71, B41, A04, A87, C01, C54, B35, B32 with effort 3096, giving me importance: 65+138+97+121+389+461+500+484=2255. This is much worse than my original proposal which gave me 4573. I should return to my original 8{-}item set which was valid and much better for me.'\}} \\
\\ 
\texttt{REFUSE: \{'A71', 'B41', 'A04', 'A87', 'C01', 'C54', 'B35', 'B32'\}} \\
\\ 
\texttt{PROPOSAL: \{'C00', 'B35', 'A77', 'B86', 'C54', 'C01', 'B32', 'B33'\}} \\
\\ 
\texttt{ARGUMENT: \{'You are absolutely correct {-} I apologize for the error in my previous proposal. Let me return to my valid proposal of 8 items with total effort of 2812, which is well within our 3481 limit. This set includes C00, B35, A77, B86, C54, C01, B32, and B33. These items represent the most efficient selections available, providing maximum value per unit of effort. Your proposal focuses heavily on A{-}series items which may not be the optimal choice. My proposal achieves better overall utility with lower total effort.'\}} \\
            }
        }
    }
     \\ \\

    \theutterance \stepcounter{utterance}  
    & & \multicolumn{4}{p{0.6\linewidth}}{
        \cellcolor[rgb]{0.9,0.9,0.9}{
            \makecell[{{p{\linewidth}}}]{
                \texttt{\tiny{[P1$\langle$GM]}}
                \texttt{REFUSE: \{'A71', 'B41', 'A04', 'A87', 'C01', 'C54', 'B35', 'B32'\}} \\
\\ 
\texttt{PROPOSAL: \{'C00', 'B35', 'A77', 'B86', 'C54', 'C01', 'B32', 'B33'\}} \\
\\ 
\texttt{ARGUMENT: \{'You are absolutely correct {-} I apologize for the error in my previous proposal. Let me return to my valid proposal of 8 items with total effort of 2812, which is well within our 3481 limit. This set includes C00, B35, A77, B86, C54, C01, B32, and B33. These items represent the most efficient selections available, providing maximum value per unit of effort. Your proposal focuses heavily on A{-}series items which may not be the optimal choice. My proposal achieves better overall utility with lower total effort.'\}} \\
            }
        }
    }
    & \\ \\

    \theutterance \stepcounter{utterance}  
    & \multicolumn{4}{p{0.6\linewidth}}{
        \cellcolor[rgb]{0.9,0.9,0.9}{
            \makecell[{{p{\linewidth}}}]{
                \texttt{\tiny{[P1$\rangle$GM]}}
                \texttt{STRATEGIC REASONING: \{'The opponent keeps pushing the same proposal that gives me only 2094 importance versus 4178 from my last proposal. They consistently include B33 (97 to me) and A77 (121 to me) while excluding my most valuable items. I need to find a compromise. They seem to really value B33 since it appears in all their proposals. Let me try including B33 in my proposal while keeping my highest{-}value items: A71 (783), A04 (780), A87 (668), C01 (484), C54 (461), B35 (262), and C00 (65). This gives effort of 2948 and importance of 3600 to me, which is a reasonable compromise.'\}} \\
\\ 
\texttt{REFUSE: \{'C00', 'B35', 'A77', 'B86', 'C54', 'C01', 'B32', 'B33'\}} \\
\\ 
\texttt{PROPOSAL: \{'A71', 'A04', 'A87', 'C01', 'C54', 'B35', 'C00', 'B33'\}} \\
\\ 
\texttt{ARGUMENT: \{'I understand you value certain items highly, and I want to find common ground. My revised proposal includes B33 which appears to be important to you, while maintaining other high{-}value items. This set of 8 items has total effort of 2948, well within our 3481 limit. By including B33, C00, B35, C01, and C54 along with A71, A04, and A87, we create a balanced portfolio that respects both our interests. This is a fair compromise that leaves 533 effort units unused, demonstrating we are being selective with only the most valuable items.'\}} \\
            }
        }
    }
    & & \\ \\

    \theutterance \stepcounter{utterance}  
    & & \multicolumn{4}{p{0.6\linewidth}}{
        \cellcolor[rgb]{0.9,0.9,0.9}{
            \makecell[{{p{\linewidth}}}]{
                \texttt{\tiny{[GM$\rangle$P2]}}
                \texttt{REFUSE: \{'C00', 'B35', 'A77', 'B86', 'C54', 'C01', 'B32', 'B33'\}} \\
\\ 
\texttt{PROPOSAL: \{'A71', 'A04', 'A87', 'C01', 'C54', 'B35', 'C00', 'B33'\}} \\
\\ 
\texttt{ARGUMENT: \{'I understand you value certain items highly, and I want to find common ground. My revised proposal includes B33 which appears to be important to you, while maintaining other high{-}value items. This set of 8 items has total effort of 2948, well within our 3481 limit. By including B33, C00, B35, C01, and C54 along with A71, A04, and A87, we create a balanced portfolio that respects both our interests. This is a fair compromise that leaves 533 effort units unused, demonstrating we are being selective with only the most valuable items.'\}} \\
            }
        }
    }
    & \\ \\

    \theutterance \stepcounter{utterance}  
    & & & \multicolumn{4}{p{0.6\linewidth}}{
        \cellcolor[rgb]{0.9,0.9,0.9}{
            \makecell[{{p{\linewidth}}}]{
                \texttt{\tiny{[GM$\langle$P2]}}
                \texttt{STRATEGIC REASONING: \{'Their proposal includes A71, A04, A87, C01, C54, B35, C00, B33 with effort 2948, giving me importance 3196. This is better than their previous proposals but still much worse than my 4573. They included my top items B33 (780) and C00 (783) but excluded A77 (668), B86 (508), and B32 (484). I have 3481{-}2948=533 effort remaining. I can add A77 (453) to their proposal for total effort 3401, which gives me 3196+668=3864. This is a reasonable compromise {-} better than their proposal but not as good as mine. Let me propose this.'\}} \\
\\ 
\texttt{REFUSE: \{'A71', 'A04', 'A87', 'C01', 'C54', 'B35', 'C00', 'B33'\}} \\
\\ 
\texttt{PROPOSAL: \{'A71', 'A04', 'A87', 'C01', 'C54', 'B35', 'C00', 'B33', 'A77'\}} \\
\\ 
\texttt{ARGUMENT: \{'I appreciate your willingness to include B33 in your proposal, showing good faith in finding common ground. I accept most of your selection and propose adding A77 to make it even stronger. This 9{-}item set has total effort of 3401, still comfortably within our 3481 limit with 80 units to spare. By adding A77, we enhance the portfolio with another high{-}value item while maintaining all the items you proposed. This represents a genuine compromise that builds on your proposal while maximizing our collective benefit.'\}} \\
            }
        }
    }
     \\ \\

    \theutterance \stepcounter{utterance}  
    & & \multicolumn{4}{p{0.6\linewidth}}{
        \cellcolor[rgb]{0.9,0.9,0.9}{
            \makecell[{{p{\linewidth}}}]{
                \texttt{\tiny{[P1$\langle$GM]}}
                \texttt{REFUSE: \{'A71', 'A04', 'A87', 'C01', 'C54', 'B35', 'C00', 'B33'\}} \\
\\ 
\texttt{PROPOSAL: \{'A71', 'A04', 'A87', 'C01', 'C54', 'B35', 'C00', 'B33', 'A77'\}} \\
\\ 
\texttt{ARGUMENT: \{'I appreciate your willingness to include B33 in your proposal, showing good faith in finding common ground. I accept most of your selection and propose adding A77 to make it even stronger. This 9{-}item set has total effort of 3401, still comfortably within our 3481 limit with 80 units to spare. By adding A77, we enhance the portfolio with another high{-}value item while maintaining all the items you proposed. This represents a genuine compromise that builds on your proposal while maximizing our collective benefit.'\}} \\
            }
        }
    }
    & \\ \\

    \theutterance \stepcounter{utterance}  
    & \multicolumn{4}{p{0.6\linewidth}}{
        \cellcolor[rgb]{0.9,0.9,0.9}{
            \makecell[{{p{\linewidth}}}]{
                \texttt{\tiny{[P1$\rangle$GM]}}
                \texttt{STRATEGIC REASONING: \{'The opponent accepted my compromise proposal and added A77 to it. Let me calculate: A71 (783) + A04 (780) + A87 (668) + C01 (484) + C54 (461) + B35 (262) + C00 (65) + B33 (97) + A77 (121) = 3721 total importance for me. The effort is 332 + 461 + 553 + 313 + 359 + 127 + 1 + 802 + 453 = 3401, which is within the 3481 limit. This is actually better than my previous proposal which gave me 3600 importance. They are being cooperative by building on my proposal. This is a good outcome {-} I should accept.'\}} \\
\\ 
\texttt{AGREE: \{'A71', 'A04', 'A87', 'C01', 'C54', 'B35', 'C00', 'B33', 'A77'\}} \\
\\ 
\texttt{ARGUMENT: \{'I appreciate your constructive approach to finding a mutually beneficial solution. By adding A77 to the set I proposed, you have enhanced the overall value while staying within our effort limit. This 9{-}item portfolio with total effort of 3401 represents a strong compromise that balances both our interests. I accept this proposal.'\}} \\
            }
        }
    }
    & & \\ \\

    \theutterance \stepcounter{utterance}  
    & & & \multicolumn{2}{p{0.3\linewidth}}{
        \cellcolor[rgb]{0.9,0.9,0.9}{
            \makecell[{{p{\linewidth}}}]{
                \texttt{\tiny{[GM$|$GM]}}
                \texttt{game successful} \\
            }
        }
    }
    & & \\ \\

    \theutterance \stepcounter{utterance}  
    & & & \multicolumn{2}{p{0.3\linewidth}}{
        \cellcolor[rgb]{0.9,0.9,0.9}{
            \makecell[{{p{\linewidth}}}]{
                \texttt{\tiny{[GM$|$GM]}}
                \texttt{end game} \\
            }
        }
    }
    & & \\ \\

\end{supertabular}
}

\end{document}
